\section{Código en Python}
    \label{anexo_python}
    Las fases 0 a 4 de la Tabla \ref{workflow} se implementan en Python y se
    reproducen completas a continuación, en su orden de ejecución.

    \subsection{\textit{dataset.py}}
        Regenera el conjunto de datos con la semilla fija 42 y añade la variable
        derivada \texttt{tarifa\_cop\_kwh} de la Tabla \ref{dataset_variables}, que
        después sirve como patrón externo contra el cual contrastar las pendientes
        estimadas.

        \vspace{10pt}

        \lstinputlisting[style=python]{utils/codes/python/dataset.py}

    \subsection{\textit{simple\_regression.py}}
        Calcula las correlaciones de la Tabla \ref{correlations}, estima el modelo M1
        de la Tabla \ref{coef_m1}, aplica las cuatro pruebas de supuestos de la Tabla
        \ref{diagnostics_m1}, cuantifica el sesgo por sector de la Tabla
        \ref{sector_bias} y produce las Figuras \ref{fit_m1}, \ref{diag_panel} y
        \ref{bias_boxplot}.

        \vspace{10pt}

        \lstinputlisting[style=python]{utils/codes/python/simple_regression.py}

    \subsection{\textit{multiple\_regression.py}}
        Estima los modelos M2 y M3, ejecuta el contraste F de la Tabla \ref{anova},
        traduce las pendientes a las tarifas de la Tabla \ref{tariffs}, compara los
        tres modelos en la Tabla \ref{model_comparison} y genera las Figuras
        \ref{fit_m3}, \ref{coef_plot} y \ref{criteria}.

        \vspace{10pt}

        \lstinputlisting[style=python]{utils/codes/python/multiple_regression.py}

    \subsection{\textit{ml\_regression.py}}
        Construye la partición estratificada y la validación cruzada de diez pliegues
        dentro de un \texttt{Pipeline}, produciendo la Tabla \ref{ml_validation} y la
        Figura \ref{holdout}.

        \vspace{10pt}

        \lstinputlisting[style=python]{utils/codes/python/ml_regression.py}

    \subsection{\textit{advanced\_viz.py}}
        Lee las tablas que dejó la fase 2 en lugar de recalcularlas y construye las
        cuatro figuras de \textit{seaborn} (\ref{heatmap}, \ref{scatter_matrix},
        \ref{lmplot} y \ref{lowess}) y las dos piezas interactivas de \textit{Plotly}
        (\ref{dashboard} y \ref{scatter_plotly}).

        \vspace{10pt}

        \lstinputlisting[style=python]{utils/codes/python/advanced_viz.py}
